\section*{Chapter 11 --- Matrices and Graphs}

\begin{solution}{exo:g12:matrix:1}
\[
A + B = \begin{pmatrix} 3 & 2\\ 1 & 2\end{pmatrix}, \quad
AB = \begin{pmatrix} 4 & 2\\ 1 & 1\end{pmatrix}, \quad
BA = \begin{pmatrix} 2 & 4\\ 1 & 3\end{pmatrix}, \quad
A^2 = \begin{pmatrix} 1 & 4\\ 0 & 1\end{pmatrix}.
\]
Note $AB \neq BA$.
\end{solution}

\begin{solution}{exo:g12:matrix:2}
$\det A = 6 - 5 = 1 \neq 0$:
$A^{-1} = \begin{pmatrix} 3 & -5\\ -1 & 2 \end{pmatrix}$.
$\det B = 12 - 12 = 0$: $B$ is not invertible.
\end{solution}

\begin{solution}{exo:g12:matrix:3}
The system is $AX = Y$ with $A$ as in \cref{exo:g12:matrix:2} and
$Y = \begin{pmatrix} 1\\ 2\end{pmatrix}$:
\[
X = A^{-1}Y = \begin{pmatrix} 3 & -5\\ -1 & 2\end{pmatrix}
\begin{pmatrix} 1\\ 2\end{pmatrix}
= \begin{pmatrix} -7\\ 3\end{pmatrix}:
\qquad x = -7,\ y = 3 .
\]
\end{solution}

\begin{solution}{exo:g12:matrix:4}
\emph{1.} $N^2 = \begin{pmatrix} 0&1\\0&0\end{pmatrix}
\begin{pmatrix} 0&1\\0&0\end{pmatrix} = 0$, and $I$ commutes with every
matrix.

\emph{2.} Since the two terms commute, the binomial theorem applies and all
terms containing $N^2$ vanish:
\[
A^k = (2I + N)^k = 2^k I + k\,2^{k-1} N
= \begin{pmatrix} 2^k & k\,2^{k-1}\\ 0 & 2^k \end{pmatrix}.
\]
Check for $k = 2$:
$A^2 = \begin{pmatrix} 2&1\\0&2\end{pmatrix}^2
= \begin{pmatrix} 4&4\\0&4\end{pmatrix}$, and the formula gives
$2^2 = 4$, $2 \times 2 = 4$. \checkmark
\end{solution}

\begin{solution}{exo:g12:matrix:5}
\emph{Induction.} For $n = 1$:
$A^1 = \begin{pmatrix} 0&1\\1&1\end{pmatrix}
= \begin{pmatrix} F_0 & F_1\\ F_1 & F_2\end{pmatrix}$. Assume the formula
for $n$; then
\[
A^{n+1} = A^n A
= \begin{pmatrix} F_{n-1} & F_n\\ F_n & F_{n+1}\end{pmatrix}
\begin{pmatrix} 0 & 1\\ 1 & 1\end{pmatrix}
= \begin{pmatrix} F_n & F_{n-1} + F_n\\ F_{n+1} & F_n + F_{n+1}\end{pmatrix}
= \begin{pmatrix} F_n & F_{n+1}\\ F_{n+1} & F_{n+2}\end{pmatrix}.
\]
\emph{Identity.} For $2\times2$ matrices, expanding shows
$\det(MN) = \det M \det N$; hence
$\det(A^n) = (\det A)^n = (-1)^n$, and
$\det A^n = F_{n-1}F_{n+1} - F_n^2$. (This is \emph{Cassini's identity}.)
\end{solution}

\begin{solution}{exo:g12:matrix:6}
\emph{1.} Ordering vertices $1, 2, 3$:
\[
M = \begin{pmatrix} 0&1&1\\ 0&0&1\\ 1&0&0\end{pmatrix}, \quad
M^2 = \begin{pmatrix} 1&0&1\\ 1&0&0\\ 0&1&1\end{pmatrix}, \quad
M^3 = \begin{pmatrix} 1&1&1\\ 1&0&1\\ 1&0&1 \end{pmatrix}.
\]

\emph{2.} $\bigl(M^3\bigr)_{11} = 1$: exactly one closed walk of length $3$
at vertex $1$, namely $1 \to 2 \to 3 \to 1$. (The walk
$1 \to 3 \to 1$ has length $2$ only, and $1 \to 3$ then $3\to1$ then
$1\to3$ ends at $3$.)
\end{solution}

\begin{solution}{exo:g12:matrix:7}
\emph{1.} $a_{n+1} = 0.8a_n + 0.3b_n$, $b_{n+1} = 0.2a_n + 0.7b_n$:
$M = \begin{pmatrix} 0.8 & 0.3\\ 0.2 & 0.7\end{pmatrix}$.

\emph{2.} $MX = X$ gives $0.8a + 0.3b = a$, \emph{i.e.} $0.3b = 0.2a$, so
$b = \frac23 a$; with $a + b = 1$: $a = 0.6$, $b = 0.4$.

\emph{3.} Using $b_n = 1 - a_n$:
$a_{n+1} = 0.8a_n + 0.3(1 - a_n) = 0.5a_n + 0.3$, so
\[
c_{n+1} = a_{n+1} - 0.6 = 0.5a_n + 0.3 - 0.6 = 0.5(a_n - 0.6) = 0.5\,c_n .
\]
Hence $c_n = 0.5^n c_0 \to 0$: $a_n \to 0.6$ and $b_n \to 0.4$, whatever the
initial distribution.
\end{solution}

\begin{solution}{exo:g12:matrix:8}
\emph{1.} $\det P = -2$, so
$P^{-1} = -\frac12\begin{pmatrix} -1 & -1\\ -1 & 1\end{pmatrix}
= \frac12\begin{pmatrix} 1 & 1\\ 1 & -1\end{pmatrix}$. Then
\[
AP = \begin{pmatrix} 4 & 2\\ 4 & -2 \end{pmatrix}, \qquad
D = P^{-1}AP = \frac12\begin{pmatrix} 1&1\\1&-1\end{pmatrix}
\begin{pmatrix} 4&2\\4&-2\end{pmatrix}
= \begin{pmatrix} 4 & 0\\ 0 & 2\end{pmatrix}.
\]

\emph{2.} From $A = PDP^{-1}$, an immediate induction gives
$A^n = PD^nP^{-1}$ with
$D^n = \begin{pmatrix} 4^n & 0\\ 0 & 2^n\end{pmatrix}$, so
\[
A^n = P D^n P^{-1}
= \begin{pmatrix} 4^n & 2^n\\ 4^n & -2^n\end{pmatrix}\cdot
\frac12\begin{pmatrix} 1&1\\1&-1\end{pmatrix}
= \frac12\begin{pmatrix} 4^n + 2^n & 4^n - 2^n\\
4^n - 2^n & 4^n + 2^n\end{pmatrix}.
\]
Applying $A^n$ to $X_0 = \begin{pmatrix} u_0\\v_0\end{pmatrix}$ reproduces
exactly the formulas of \cref{ex:g12:matrix:coupled}.
\end{solution}
